\section{Data Visualization}
\label{sec:datavis}

Visualization data transforms it into information. Data itself is just noise, whereas information holds value. In what way data needs to be transformed is determined by context [\cite{terganKnowledgeInformationVisualization2005}].

The type of visualization depends on the purpose of the information. The main constraint in this case is cognitively demanding situations.

\subsection{Situation Awareness}

As mentioned earlier, the tasks to execute are mentally demanding. Concentrating on fulfilling task objectives while maintaining awareness of the surrounding area result in a high cognitive load, while associated risks prohibit human error. Endsley's model of situation awareness can be referenced to estimate the requirements of information acquisition in dynamic situations.

Situation awareness is a model to describe the process of decision-making in dynamic situations.

Displaying relevant information at a glance can therefore improve SA directly. Aside from that it will free up cognitive load, which can then be reallocated.

By displaying certain information [like oxygen etc] on a HUD, at least some level of situational awareness could be restored. Going by Endsley's definition, that would increase the likelihood of making better decisions.

ENDSLEY 1995
ENDSLEY 1995 proposes a model of situation awareness which includes 3 levels, namely perception, comprehension and projection. Only Level 1 awareness—perception—is relevant for UI design itself, though some concepts from Level 2 can be used to exemplify things to avoid.

\subsection{Types of Information}

Depending on the character of the data, different methods of visualization could be chosen.

Comparable to basic dashboard data, only value of the parameter is relevant, since there is no relationship between the data. At most it is the relation between a current and a past value of a parameter.

Purposes of color to cenvey information are grouping or relationships. Grouping is necessary for the HUD, but is managed by the spatial location of and the spacing between elements.

WILKE2019 fundamentals dataviz ch2
https://learning.oreilly.com/library/view/fundamentals-of-data/9781492031079/ch02.html

\subsection{Visualization Methods}